\documentclass{article}
\usepackage{mathnotes}

\title{Projection}
\description{Vector projection formulas and geometric interpretation, including component calculation and orthogonal projection methods.}

\begin{document}

\section{Projection}

\subsection{Definition of Projection}


\begin{definition}[projection]\synonyms{vector projection", "orthogonal projection"}
	Take two \@{vectors} in $R^n$, $\vec{u}$ and $\vec{v}$, with $\vec{v} \neq \vec{0}$.

	The \term[scalar component]{scalar projection} of $\vec{u}$ onto $\vec{b}$ is a scalar equal to $a_1 = \norm{\vec{a}} \cos{\theta},$ where $\theta$ is the \@{angle} between $\vec{a}$ and $\vec{b}.$ It is denoted as follows, with an equivalent formula:

	\[ \comp_{\vec{v}} \vec{u} =  \norm{\proj_{\vec{v}} \vec{u}} = \frac{\inner{\vec{u}}{\vec{v}}}{\norm{\vec{v}}} \]


	The \textbf{projection} of $\vec{u}$ onto $\vec{v}$ is the vector whose \@{magnitude} is the \@{scalar projection} of $\vec{u}$ onto $\vec{v}$ with the same \@{direction} as $\vec{v}$:

	\[ \proj_{\vec{v}} \vec{u} = \frac{\inner{\vec{u}}{\vec{v}}}{\norm{\vec{v}}^2} \vec{v} = (\norm{\vec{a}}\cos{\theta})\vec{\hat{b}}. \]

	The \term[vector rejection]{rejection} of $\vec{u}$ from $\vec{v}$ is the component of $\vec{u}$ that is \@{orthogonal} to $\vec{v}$ and is given as

        \[ \oproj_\vec{v} \vec{u} = \vec{u} -  \proj_{\vec{v}} \vec{u}. \]
\end{definition}

\begin{intuition}
One way to visualize this is to shine a flashlight from a direction \@{perpendicular} to $\vec{v}$, with $\vec{u}$ between the flashlight and $\vec{v}$ (pretend the vectors are opaque). The shadow the flashlight casts onto $\vec{v}$ will be the projection of $\vec{u}$ onto $\vec{v}$.
\end{intuition}

\begin{theorem}[Orthogonal Decomposition]
	For $\vec{a}, \vec{b} \in \reals^n$, with $\vec{b} \neq \vec{0},$ 

	\[ \vec{a} = \proj_{b} a + \oproj_{b} a. \]
\end{theorem}


\begin{theorem}[Pythagorean Theorem]
	For \@{orthogonal} \@{vectors} $\vec{a}, \vec{b} \in \reals^n,$

	\[ \norm{\vec{a}}^2 + \norm{\vec{b}}^2 = \norm{\vec{a} + \vec{b}}^2. \]
\end{theorem}
\begin{proof}
	By definition of the \@{euclidean norm},
	\begin{align*}
	\norm{\vec{a} + \vec{b}}^2 & = \inner{\vec{a} + \vec{b}}{\vec{a} + \vec{b}} \\
	                          & = \norm{\vec{a}}^2 + 2 \inner{\vec{a}}{\vec{b}} + \norm{\vec{b}}^2.
	\intertext{Because $\vec{a}$ and $\vec{b}$ are \@{orthogonal}, their \@{euclidean inner product} is $\vec{0},$ and we have}
	\norm{\vec{a} + \vec{b}}^2 & = \norm{\vec{a}}^2 + \norm{\vec{b}}^2.
	\end{align*}

\end{proof}

\begin{theorem}[Cauchy-Schwarz Inequality]
For any \@{vectors} $\vec{x},\vec{y} \in \reals^n,$

	\[ \abs{\inner{\vec{x}}{\vec{y}}} \leq \norm{\vec{x}} \norm{\vec{y}}, \]

	where $\inner{\vec{x}}{\vec{y}}$ denotes the \@{euclidean inner product} and $\norm{\vec{x}} = \sqrt{ \inner{\vec{x}}{\vec{x}} }$ is the \@{euclidean norm}.
\end{theorem}
\begin{proof}
From \@{orthogonal decomposition} we have
	\begin{align*}
		\vec{a} & = \proj_{b} a + \oproj_{b} a.
		\shortintertext{Taking the \@{euclidean norm} squared of both sides gives }
		\norm{\vec{a}}^2 & = \norm{\proj_{b} a + \oproj_{b} a}^2 \\
		              & = \norm{\proj_{b} a}^2 + 2 \inner{\proj_{b} a}{\oproj_{b} a} + \norm{\oproj_{b} a}^2 \\
		              & = \norm{\proj_{b} a}^2 + \norm{\oproj_{b} a}^2, \\
		\shortintertext{with the middle term dropping out due to orthogonality. Now, because $\norm{\oproj_{b} a}^2 \geq 0$,}
		\norm{\vec{a}}^2 & \geq \norm{\proj_{b} a}^2. \\
		\shortintertext{Now we'll flip sides and expand \@{projection} according to its definition}
		 \norm{\proj_{b} a}^2 & \leq \norm{\vec{a}}^2  \\
		 \norm{\frac{ \inner{\vec{a}}{\vec{b}} }{\norm{\vec{b}}^2} \vec{b}}^2 & \leq \\
		 \parens{\frac{ \inner{\vec{a}}{\vec{b}} }{\norm{\vec{b}}^2}}^2 \norm{\vec{b}}^2 & \leq \\
		 \frac{ \inner{\vec{a}}{\vec{b}}^2 }{\norm{\vec{b}}^2} & \leq \norm{\vec{a}}^2 \\
		 \inner{\vec{a}}{\vec{b}}^2 & \leq \norm{\vec{a}}^2 \norm{\vec{b}}^2. \\
		 \shortintertext{Taking the square root of both sides we get what we want to show:}
		 \sqrt{\inner{\vec{a}}{\vec{b}}^2} & \leq \sqrt{\norm{\vec{a}}^2 \norm{\vec{b}}^2} \\
		 \inner{\vec{a}}{\vec{b}} & \leq \norm{\vec{a}} \norm{\vec{b}}.
	\end{align*}
\end{proof}
\begin{intuition}
	This is saying that the length of the \@{projection} of $\vec{a}$ onto $\vec{b}$ cannot exceed the length of $\vec{a}.$
\end{intuition}


\end{document}
